\section{Introduction}
\label{sec:intro}

A person who communicates through an AAC board composes a message by selecting symbols, and a speech-generating device voices the selected labels.
The result is telegraphic: \emph{Water}, or \emph{No}+\emph{Pizza}. A language model can expand such selections into sentences, and Indian-language boards would benefit from this because Hindi, Gujarati and Marathi require agreement, postpositions and negation forms that a literal join of labels cannot supply.
For AAC, however, fluency is not sufficient, because the device speaks on behalf of the user and the listener cannot see the selection.
Three outputs of the cascade evaluated in this paper illustrate the risk: for \emph{Water} it produced ``I need food.'' (a concept replaced); for \emph{No}+\emph{Water} in Marathi it produced \hi{मला पाणी हवे आहे.} (``I want water''; polarity inverted); and for \emph{Tired} it produced ``The user is tired of waiting for their meal and wants to go out to eat.'' (content added).

\sys{} addresses this with a \emph{concept gate}. The selection is converted into a ledger of required and forbidden concepts, a prompted cascade realizes a sentence, a deterministic validator accepts it only if the ledger is satisfied, and a template built from the user's own labels replaces it otherwise.
The system adds a four-locale symbol board, trained next-word predictors and zero-shot voice synthesis. We contribute:
\begin{itemize}
\item the gated realization pipeline and its acceptance rule (\S\ref{sec:system}), with an API that carries polarity and intent explicitly;
\item a frozen, paired evaluation with baselines, an independent scorer and a failure analysis (\S\ref{sec:eval}), and the fixtures and scorer for reuse;
\item trained English and Hindi predictors, evaluated on held-out conversations, and the end-to-end system with its demonstration walkthrough (Figures~\ref{fig:arch} and~\ref{fig:ui}).
\end{itemize}
